\documentclass{subfiles}
\begin{document}
\chapter{Methods}
This chapter will present the methods behind all our results, both intermediate analyses and the final results. We will attempt to have some form of chronological order in the presentation of the methods, so the reader can follow the development of the project. Meaning, we will start by presenting the methods from the background analyis done on the system, and then move onwards to the methods that consitute our main goal of the project. As we are working with a rather unknown double well potential, expressed (hopefully) as a distinguishable system, there are many aspects we should study to confirm the viability of our approach. All of our methods are implemented numerically in Python, and some of the methods will be presented in a general sense, whilst more technical details will be presented as pseudo-code. All code used in the project is readily available at: \url{https://github.com/Jonnyigeh/Thesis} \textcolor{red}{(TODO: Make repo public, and CLEAN IT)}.
\begin{figure}
\centering
\includegraphics[width=0.5\textwidth]{figs/placeholder_image.jpeg}
\caption{TODO: Replace this placeholder image with a relevant figure. }
\label{fig:methods_intro}
\end{figure}
\subfile{../sections/chap2/hartree.tex}
\subfile{../sections/chap2/general_study.tex}
\subfile{../sections/chap2/time_evolution_methods.tex}
\subfile{../sections/chap2/summary_workflow.tex}
\end{document}